\section{Duality and the monopole conjectures}\label{sec:sduality}

\subsection{The $SL(2,\ZZ)$ conjecture}\label{ssec:sl2z}

The conjecture of \cite{Sen94duality} (claim 1 of Table~\ref{tab:catalogue}) states that the heterotic string compactified on $T^6$, whose low energy theory is $\mathcal N=4$ supergravity coupled to 22 vector multiplets, is invariant under the group $SL(2,\ZZ)$ generated by the shift of the axion and by the inversion $\lambda\mapsto -1/\lambda$ of the complex coupling $\lambda = a + i e^{-\phi}$, acting on the electric and magnetic charge vectors $(Q,P)$ of the $O(6,22)$ lattice by the same matrix.\footnote{The conjecture was made in two steps. Schwarz and Sen \cite{SchwarzSen93} showed that the equations of motion of the low energy theory, not its action, are invariant under $SL(2,\RR)$, and that the charge lattice reduces the symmetry to $SL(2,\ZZ)$; \cite{Sen94duality} then proposed that the full string theory, at every order and beyond perturbation theory, has this symmetry. The $SL(2,\RR)$ of the equations of motion is a theorem of the supergravity theory; it is the lift to the full string theory that is the conjecture.} The inversion exchanges weak and strong coupling, so no perturbative computation can test it directly. Sen reduced the test to the spectrum of states that saturate the Bogomolnyi bound, whose masses and degeneracies are protected by supersymmetry: the conjecture requires that for every pair of coprime integers $(p,q)$ there exist states carrying $p$ units of electric and $q$ units of magnetic charge, in the same supermultiplet as the elementary string states with $(1,0)$.

For $q=1$ these are the dyons of the $\mathcal N=4$ theory and they were known to exist \cite{Osborn79}. For $(p,q) = (1,2)$ the conjecture requires a bound state of two monopoles and one unit of electric charge, and nothing in the classical theory produces such a state: the static BPS monopoles have no net force, but a normalisable quantum bound state is not supplied by a static separated configuration and there is no classical bound state. The prediction is therefore a genuine one, and it was tested in \cite{Sen94dyon}, as described in the next subsection. The subsequent tests of the conjecture are the three-string junction states \cite{Sen98network,Sen07junction}, the counting of quarter-BPS dyons (Section~\ref{sec:dyons}), whose partition function is a Siegel form invariant under the duality group, the match of the black hole entropy with the index (Section~\ref{sec:entropy}), and the equivalent statement for $\mathcal N=4$ super-Yang--Mills theory, where the conjecture of Montonen and Olive \cite{MontonenOlive77} is tested by the same bound states \cite{Sen94dyon,Porrati96,GauntlettLowe96,LeeWeinbergYi96}, by the modular properties of partition functions on four-manifolds \cite{VafaWitten94}, and by its embedding into the geometric Langlands programme \cite{KapustinWitten07}.

\begin{verdict}[claim 1]\label{ver:sl2z}
Confirmed. Every protected quantity that has been computed is consistent with the $SL(2,\ZZ)$ symmetry, and the strongest test, the exact dyon index of Section~\ref{sec:dyons}, is a Siegel form of the duality group. \gap{A proof of the invariance beyond the protected quantities is not known.} It would require a non-perturbative definition of the heterotic string on $T^6$, which has no definition at strong coupling other than through such dualities, and the tests cited here do not constitute such a construction.
\end{verdict}

One point about the statement itself should be made precisely, because it is the kind of place where a conjecture can be ``wrong'' without its author being wrong. For gauge theories with a non-simply-laced gauge group the duality group of the $\mathcal N=4$ theory is not $SL(2,\ZZ)$: the inversion maps the gauge group $G$ to its Langlands dual $G^\vee$, and the group generated by the shift and the inversion is a Hecke group $\Gamma_0(q)$-type subgroup, with $q=2$ or $3$ \cite{GGPZ95,DFHK96,AKS06}. This is not a correction of \cite{Sen94duality} or of \cite{Sen94dyon}: both papers treat $SU(2)$ and the heterotic string, where the duality group is $SL(2,\ZZ)$, and \cite{Sen94dyon} is explicit that the test is for $SU(2)$. We mention it because the literature sometimes attributes an $SL(2,\ZZ)$ conjecture ``for all groups'' to the 1994 papers; they contain no such statement.

\subsection{The two-monopole bound state}\label{ssec:twomonopole}

The moduli space of two $SU(2)$ monopoles of total charge two is $\RR^3\times (S^1\times \mathcal M_0)/\ZZ_2$, where $\mathcal M_0$ is the Atiyah--Hitchin manifold \cite{AtiyahHitchin88}, a complete hyperk\"ahler four-manifold with an $SO(3)$ action and a metric known in closed form (for the hypercomplex structures underlying such metrics see also \cite{Bha22b}). The low energy dynamics of two monopoles is the supersymmetric quantum mechanics on this moduli space, and a bound state of two monopoles with one unit of electric charge and the right supermultiplet structure exists if and only if the Atiyah--Hitchin manifold carries a harmonic two-form that is square integrable, anti-self-dual (with the orientation fixed by the complex structures), odd under the $\ZZ_2$, and unique.\footnote{The electric charge of the state is the momentum along the $S^1$ factor, and the $\ZZ_2$ acts on it; the state with one unit of electric charge requires the form on $\mathcal M_0$ to be odd under the $\ZZ_2$ which combines with the half shift on $S^1$. The multiplet structure is fixed by the fermion zero modes: the state must be a single short multiplet of $\mathcal N=4$, which is the statement that the harmonic form is unique and of a definite duality type. Sen's paper carries this out in detail and we do not repeat it.} Sen wrote down the form. In the coordinates of \cite{AtiyahHitchin88}, with the metric $ds^2 = f^2 d\rho^2 + a^2\sigma_1^2 + b^2\sigma_2^2 + c^2\sigma_3^2$ on left-invariant forms with $d\sigma_i = \tfrac12\epsilon_{ijk}\sigma_j\wedge\sigma_k$, the form is
\begin{equation}\label{eq:senform}
\begin{aligned}
\omega &= F(\rho)\,\Big(d\sigma_1 - \frac{fa}{bc}\, d\rho\wedge\sigma_1\Big),\\
F(\rho) &= F_0\exp\Big(-\int_{\pi}^{\rho}\frac{fa}{bc}\,d\rho'\Big),
\end{aligned}
\end{equation}
which is anti-self-dual by the choice of the coefficient of $d\rho\wedge\sigma_1$, closed because $dF/d\rho = -(fa/bc)F$, and square integrable because $fa/(bc)\to\tfrac12$ as $\rho\to\infty$ (where $f\simeq -1$, $a\simeq b\simeq\rho$, $c\simeq -2$) makes $F$ decay as $e^{-\rho/2}$ while the volume grows polynomially; near the bolt $\rho\to\pi$, in the coordinates adapted to it, the form is smooth.\footnote{Equations (16) to (19) of \cite{Sen94dyon}. The form is invariant under the identification $(\rho,\theta,\phi,\psi)\sim(\rho,\pi-\theta,\pi+\phi,-\psi)$ that defines the Atiyah--Hitchin manifold and changes sign under $\psi\mapsto\psi+\pi$, which is the $\ZZ_2$ of the two-monopole moduli space; this is the parity required for odd electric charge. The uniqueness of the form among $L^2$ harmonic two-forms was argued in \cite{Sen94dyon} from the symmetries; the complete statement that the space of $L^2$ harmonic forms of $\mathcal M_0$ is one-dimensional is Hitchin's theorem below. Sen notes that the same form had appeared in work of Gibbons and Ruback for another purpose.}

\begin{theorem}[Sen 1994; Hitchin 2000]\label{thm:atiyahhitchin}
The space of $L^2$ harmonic forms on the Atiyah--Hitchin manifold is one-dimensional. It is spanned by the two-form \eqref{eq:senform}, which is anti-self-dual and odd under the $\ZZ_2$ that acts on the moduli space of two monopoles.
\end{theorem}

\begin{proof}
Existence, anti-self-duality, square integrability and the $\ZZ_2$ parity are the computation of \cite{Sen94dyon} summarised above; we have checked each step. The uniqueness result is taken from \cite{Hitchin00}. We do not infer it merely by counting a circle fixed set: orientation and the particular cover of the monopole moduli space matter.
\end{proof}

\begin{verdict}[claim 2]
Proved. The bound state predicted by $SL(2,\ZZ)$ exists, and it is the only one in the two-monopole sector, which is also what the duality predicts: the $(1,2)$ state must be a single short multiplet.
\end{verdict}

\subsection{The Sen conjecture on $L^2$ cohomology}\label{ssec:senconjecture}

The argument of \cite{Sen94dyon} applies to any charge. Let $\mathcal M_k^0$ denote the reduced moduli space of $k$ monopoles, of real dimension $4k-4$, and let $\mathcal M_k = \RR^3\times(S^1\times\mathcal M^0_k)/\ZZ_k$. The states with $k$ units of magnetic charge and $p$ units of electric charge, with $\gcd(p,k)=1$, correspond to $L^2$ harmonic forms on $\mathcal M^0_k$ transforming under the $\ZZ_k$ by $e^{2\pi i p/k}$; $S$-duality requires exactly one such form, in the middle degree $2k-2$, for every $p$ coprime to $k$, and none for other $p$ (which would give states not in the duality orbit of the elementary string states). This is the \emph{Sen conjecture}; Segal and Selby \cite{SegalSelby96} proved that the image of the compactly supported cohomology in the ordinary cohomology, which is a lower bound for the $L^2$ cohomology, has exactly the predicted dimension in the middle degree and vanishes elsewhere. The conjecture is thus reduced to the statement that every $L^2$ harmonic form comes from compactly supported cohomology.

\begin{theorem}[published partial results and their hypotheses]\label{thm:senconj}
The Segal--Selby image has the dimensions predicted by the conjecture \cite{SegalSelby96}. Hitchin's middle-degree result and the charge-two case are established in \cite{Hitchin00}. Kottke--Rochon \cite[Theorem~A]{KottkeRochon24} prove the charge-three case for a monopole metric with a quasi-fibred boundary structure; that the metric of $\mathcal M^0_3$ has this structure is a hypothesis of their theorem. These hypotheses are part of the statements, not results proved here.
\end{theorem}
\begin{proof}
The theorem records the cited results with their own hypotheses. Quoting Theorem~A of \cite{KottkeRochon24} without its metric hypothesis would overstate it.
\end{proof}
\begin{verdict}[claim 3]
Confirmed: proved for $k=2$, and for $k=3$ under the metric hypothesis of \cite{KottkeRochon24}. \gap{The statement for every $k\ge3$ and every charge sector is not proved.} The open part is a question about the $L^2$ cohomology of the monopole moduli spaces, not evidence against the conjecture.
\end{verdict}

\subsection{Duality-symmetric actions, the Kerr--Sen black hole and string networks}\label{ssec:misc}

Three further claims of this period are theorems and we record them briefly.

\begin{proposition}[claims 4, 5 and 35]\label{prop:misc}
\begin{enumerate}[leftmargin=1.6em,itemsep=1pt]
\item[(i)] Schwarz and Sen \cite{SchwarzSen94} wrote an action with two gauge potentials and a first-order Lagrangian; it is invariant under the duality rotation and has equations of motion equivalent to Maxwell's equations; the same construction gives a duality-symmetric form of the heterotic effective action in four dimensions.
\item[(ii)] The metric, dilaton, gauge field and axion of \cite{Sen92kerr} solve the equations of motion of the low energy heterotic string in four dimensions; the solution is obtained from the Kerr solution by the $O(1,1)$ twisting transformation of the string background, which maps solutions to solutions.
\item[(iii)] A planar network of $(p,q)$-strings of type IIB with the charges conserved at each junction and the angles fixed by the tensions is supersymmetric, preserving one quarter of the supersymmetry. The proof is in \cite{Sen98network} itself, from the supersymmetry algebra of the $(p,q)$-string; the three-string junction was shown to be a BPS state by \cite{ReyYee98} and, lifted to a holomorphic membrane in M-theory, by \cite{MatsuoOkuyama98}. The statement was a theorem when it was made, and the catalogue records it as proved.
\end{enumerate}
\end{proposition}

\begin{proof}
(i) is a direct computation in the cited paper: the two potentials $A^{(1)}, A^{(2)}$ enter through a first-order Lagrangian whose equations of motion are $F^{(2)} = *F^{(1)}$ up to gauge, so that eliminating $A^{(2)}$ recovers the Maxwell equations; the Lorentz invariance is not manifest but is verified. (ii) is the statement that the twisting transformation of \cite{Sen92kerr}, which acts on backgrounds with a Killing direction, is a symmetry of the equations of motion; this was proved for the full string effective action at lowest order in $\alpha'$, and the Kerr--Sen metric is the image of Kerr. (iii) follows from the supersymmetry algebra: the BPS condition for a network is linear in the charge vectors at each junction, and the tension balance at a junction is the statement that the central charges add; \cite{Sen98network} verifies this and \cite{Sen07junction} uses the same networks to count quarter-BPS dyons. For (iii), ``proved'' refers to the effective BPS configuration; existence and completeness of the quantum spectrum at finite coupling are separate physical questions.
\end{proof}
